% -------------------------------------------------------
%  Abstract
% -------------------------------------------------------
\بدون‌تورفتگی \مهم{چکیده}: در این گزارش سه موضوع مرتبط با لایه‌ی فیزیکی و معماری شبکه بررسی می‌شود.
ابتدا کابل‌های هم‌محور، زوج‌به‌هم‌تابیده و فیبر نوری از نظر سرعت، افت سیگنال، احتمال خطا و صرفه‌ی اقتصادی مقایسه می‌شوند.
سپس وظایف چهار لایه‌ی معماری \لر{TCP/IP} و تفاوت آن با مدل مرجع هفت‌لایه‌ای \لر{OSI} توضیح داده می‌شود.
در پایان، سازوکار \لر{Auto-MDI/MDIX} و علت کارکرد کابل‌های مستقیم در تجهیزات امروزی بررسی می‌گردد.

\بدون‌تورفتگی \مهم{واژه‌های کلیدی}:
رسانه‌ی انتقال، کابل هم‌محور، زوج‌به‌هم‌تابیده، فیبر نوری، \لر{TCP/IP}، \لر{OSI}، \لر{Auto-MDI/MDIX}

\section{مقدمه}

انتخاب رسانه‌ی فیزیکی بر ظرفیت، برد، پایداری و هزینه‌ی شبکه اثر مستقیم دارد.
از سوی دیگر، معماری لایه‌ای امکان می‌دهد وظایف پیچیده‌ی ارتباط شبکه به بخش‌های مشخص تقسیم شوند.
سه پرسش این آزمایش به همین دو جنبه و نیز به یکی از قابلیت‌های مهم رابط‌های اترنت جدید مربوط‌اند.
پاسخ‌ها با تکیه بر اسناد رسمی \لر{IEEE}، \لر{IETF}، \لر{ISO} و \لر{ITU-T} و مستندات فنی سازندگان ارائه شده‌اند.

\section{پاسخ پرسش‌ها}

\subsection{مقایسه‌ی رسانه‌های انتقال}

سرعت، احتمال خطا و افت سیگنال فقط به نام رسانه وابسته نیستند؛
رده و ساخت کابل، طول مسیر، فرکانس، فرستنده‌ـ‌گیرنده، کیفیت اتصال و نویز محیط نیز مؤثرند.
بنابراین مقایسه‌ی درست باید با ذکر یک استاندارد یا نمونه‌ی مشخص انجام شود.

\subsubsection{کابل هم‌محور}

کابل هم‌محور\پاورقی{Coaxial} به‌دلیل شیلد رسانای پیرامون هادی مرکزی، در برابر نویز الکترومغناطیسی از \لر{UTP} مقاوم‌تر است.
در اترنت قدیمی، \لر{10BASE2} نرخ ۱۰ مگابیت‌برثانیه را تا ۱۸۵ متر و \لر{10BASE5} همین نرخ را تا ۵۰۰ متر فراهم می‌کرد [۱].
افت کابل با فرکانس افزایش می‌یابد؛ برای نمونه، برگه‌ی مشخصات یک کابل \لر{RG-58} شرکت \لر{Belden} افت نامی ۴٫۷ دسی‌بل بر ۱۰۰ متر در ۱۰ مگاهرتز و ۲۳ دسی‌بل بر ۱۰۰ متر در ۲۰۰ مگاهرتز را اعلام می‌کند [۲].
شیلد باعث کاهش ورود نویز می‌شود، ولی اتصال نامناسب، ناهماهنگی امپدانس یا حذف مقاومت انتهایی می‌تواند بازتاب و خطا ایجاد کند.

کابل هم‌محور برای زیرساخت موجود تلویزیونی و رادیویی، دوربین مداربسته یا سامانه‌های قدیمی مقرون‌به‌صرفه است؛
اما برای یک شبکه‌ی محلی اترنت جدید، توپولوژی باس، عیب‌یابی دشوار و سرعت کم استانداردهای قدیمی آن معمولاً توجیه اقتصادی ندارد.

\subsubsection{کابل زوج‌به‌هم‌تابیده}

کابل زوج‌به‌هم‌تابیده\پاورقی{Twisted Pair} رایج‌ترین انتخاب برای کابل‌کشی افقی شبکه‌های محلی است.
تاباندن زوج‌ها نویز القایی مشترک و هم‌شنوی را کم می‌کند و نوع \لر{STP} با شیلد، مصونیت بیشتری از \لر{UTP} دارد.
استانداردهای اترنت نمونه‌هایی از ۱۰۰ مگابیت‌برثانیه روی دو زوج \لر{Cat5} تا ۱ گیگابیت‌برثانیه روی چهار زوج \لر{Cat5} در فاصله‌ی ۱۰۰ متر را تعریف کرده‌اند [۱، ۳]؛
در کابل‌کشی جدید، \لر{Cat6A} می‌تواند ۱۰ گیگابیت‌برثانیه را تا ۱۰۰ متر پشتیبانی کند [۴].

افت و هم‌شنوی با افزایش طول و فرکانس بیشتر می‌شوند و \لر{UTP} در محیط پرنویز از فیبر و کابل هم‌محور آسیب‌پذیرتر است؛
بااین‌حال، اگر کابل استاندارد و سالم باشد، سامانه برای نرخ خطای قابل‌قبول طراحی شده است.
قیمت کم کابل و سوکت \لر{RJ45}، نصب و تعمیر ساده، فراوانی تجهیزات و امکان انتقال توان \لر{PoE} سبب می‌شود این رسانه برای اتصال رایانه، تلفن، دوربین و نقطه‌ی دسترسی در فاصله‌های حداکثر ۱۰۰ متر اقتصادی‌ترین گزینه باشد.

\subsubsection{فیبر نوری}

فیبر نوری داده را به‌صورت نور انتقال می‌دهد، در برابر تداخل الکترومغناطیسی ذاتاً مصون است و به‌علت جداسازی الکتریکی، اختلاف پتانسیل زمین و صاعقه را نیز از مسیر کابل منتقل نمی‌کند.
در نتیجه، اگر کانکتورها تمیز و سالم و شعاع خمش مجاز رعایت شود، احتمال خطای ناشی از نویز محیطی از دو رسانه‌ی مسی کمتر است.
افت آن نیز در مسافت‌های بلند بسیار کمتر است؛ برای نمونه، \لر{ITU-T G.652} برای فیبر تک‌حالته حداکثر ضریب تضعیف ۰٫۴ دسی‌بل بر کیلومتر در ۱۳۱۰ نانومتر و ۰٫۳۵ دسی‌بل بر کیلومتر در ۱۵۵۰ نانومتر را مشخص می‌کند [۵].

فیبر ظرفیت بسیار بالا و بردی از صدها متر تا چندین کیلومتر دارد، ولی فرستنده‌ـ‌گیرنده‌ی نوری، جوش و ابزار آزمون آن معمولاً گران‌تر و نصبش تخصصی‌تر است [۴].
بنابراین برای ستون فقرات ساختمان یا دانشگاه، ارتباط بین ساختمان‌ها، مسیر بیش از ۱۰۰ متر، محیط صنعتی پرنویز و نیاز به پهنای باند یا توسعه‌ی آینده توجیه دارد؛
برای یک میز کار نزدیک سوئیچ، زوج‌به‌هم‌تابیده معمولاً ارزان‌تر است.

خلاصه‌ی این مقایسه در جدول~\رجوع{جدول:مقایسه‌رسانه‌ها} آمده است.

\شروع{لوح}[ht]
\تنظیم‌ازوسط
\شرح{مقایسه‌ی کیفی رسانه‌های انتقال در کاربردهای متداول شبکه}

\شروع{جدول}{|p{2.4cm}|p{3.4cm}|p{3.7cm}|p{4.2cm}|}
\خط‌پر
\سیاه رسانه & \سیاه سرعت و برد نمونه & \سیاه افت و احتمال خطا & \سیاه کاربرد اقتصادی \\
\خط‌پر \خط‌پر
هم‌محور & ۱۰ مگابیت‌برثانیه تا ۱۸۵ یا ۵۰۰ متر در اترنت قدیمی & افت متوسط؛ مقاوم‌تر از \لر{UTP} در برابر نویز، حساس به ناهماهنگی امپدانس & زیرساخت هم‌محور موجود، تلویزیون، رادیو و سامانه‌های قدیمی \\
\خط‌پر
زوج‌به‌هم‌تابیده & ۱ تا ۱۰ گیگابیت‌برثانیه تا ۱۰۰ متر، متناسب با رده & افت و هم‌شنوی بیشتر در فرکانس بالا؛ \لر{STP} مقاوم‌تر از \لر{UTP} & اتصال تجهیزات شبکه‌ی محلی، هزینه‌ی کم، نصب ساده و \لر{PoE} \\
\خط‌پر
فیبر نوری & ظرفیت بسیار بالا؛ از صدها متر تا چندین کیلومتر & کمترین افت در مسافت بلند و مصون از نویز الکترومغناطیسی & ستون فقرات، ارتباط بین ساختمان‌ها، برد یا پهنای باند زیاد و محیط پرنویز \\
\خط‌پر
\پایان{جدول}

\برچسب{جدول:مقایسه‌رسانه‌ها}
\پایان{لوح}

از نظر کیفی و در یک نصب صحیح، فیبر کمترین حساسیت را به نویز و کمترین افت را در مسافت بلند دارد؛
کابل هم‌محور به‌کمک شیلد از \لر{UTP} مقاوم‌تر است؛
و زوج‌به‌هم‌تابیده با وجود برد کوتاه‌تر، بهترین توازن هزینه و کارایی را برای شبکه‌ی محلی فراهم می‌کند.
احتمال خطای عددی را نمی‌توان فقط از نوع کابل تعیین کرد، زیرا کدگذاری خط، نسبت سیگنال‌به‌نویز، طول و کیفیت گیرنده نیز در آن دخیل‌اند.

\subsection{معماری \لر{TCP/IP} و مقایسه با \لر{OSI}}

معماری \لر{TCP/IP} مجموعه پروتکل عملی اینترنت است و معمولاً در چهار لایه بیان می‌شود.
سند \لر{RFC 1122} همین چهار لایه‌ی کاربرد، انتقال، اینترنت و پیوند را معرفی می‌کند [۶].
داده هنگام ارسال از بالا به پایین کپسوله و در مقصد به‌ترتیب معکوس باز می‌شود.

\شروع{شمارش}

\فقره
\مهم{لایه‌ی کاربرد\پاورقی{Application}:}
خدمات موردنیاز برنامه‌های شبکه را فراهم می‌کند.
پروتکل‌هایی مانند \لر{HTTP}، \لر{DNS}، \لر{SMTP} و \لر{SSH} در این سطح قرار می‌گیرند.
در \لر{TCP/IP} برای وظایف نشست و نمایش، لایه‌های مستقل اجباری وجود ندارد و برنامه یا پروتکل کاربردی قالب داده، رمزگذاری، فشرده‌سازی و مدیریت نشست را در صورت نیاز انجام می‌دهد.

\فقره
\مهم{لایه‌ی انتقال\پاورقی{Transport}:}
ارتباط انتها‌به‌انتها و چندگانه‌سازی برنامه‌ها با شماره‌ی درگاه را بر عهده دارد.
\لر{TCP} یک جریان اتصال‌گرا، مرتب و قابل‌اعتماد با کنترل جریان و ازدحام ارائه می‌کند؛
\لر{UDP} سرویس دیتاگرامی بدون تضمین تحویل و ترتیب ارائه می‌دهد و سربار کمتری دارد [۶].

\فقره
\مهم{لایه‌ی اینترنت\پاورقی{Internet}:}
بسته را میان شبکه‌های متفاوت از مبدأ به مقصد می‌رساند.
\لر{IP} آدرس‌دهی منطقی و ارسال دیتاگرام را فراهم می‌کند و روترها بر پایه‌ی نشانی مقصد مسیر بعدی را انتخاب می‌کنند.
این لایه سرویس بهترین‌تلاش و بدون تضمین تحویل ارائه می‌دهد؛
\لر{ICMP} نیز گزارش خطا و پیام‌های کنترلی را پشتیبانی می‌کند [۶].

\فقره
\مهم{لایه‌ی پیوند یا دسترسی شبکه\پاورقی{Link/Network Access}:}
انتقال روی شبکه‌ی مستقیماً متصل را انجام می‌دهد و جزئیاتی مانند قاب‌بندی، آدرس \لر{MAC}، دسترسی به رسانه، تشخیص خطای قاب و ارسال بیت روی مس یا فیبر را پوشش می‌دهد.
\لر{Ethernet} و \لر{Wi-Fi} نمونه‌های این لایه‌اند.

\پایان{شمارش}

نگاشت تقریبی لایه‌های دو معماری در جدول~\رجوع{جدول:نگاشت‌لایه‌ها} نشان داده شده است.

\شروع{لوح}[ht]
\تنظیم‌ازوسط
\شرح{نگاشت تقریبی لایه‌های \لر{OSI} و \لر{TCP/IP}}

\شروع{جدول}{|c|c|c|}
\خط‌پر
\سیاه شماره‌ی لایه‌ی \لر{OSI} & \سیاه لایه‌ی \لر{OSI} & \سیاه لایه‌ی متناظر \لر{TCP/IP} \\
\خط‌پر \خط‌پر
۷ & کاربرد & کاربرد \\
۶ & نمایش & کاربرد \\
۵ & نشست & کاربرد \\
۴ & انتقال & انتقال \\
۳ & شبکه & اینترنت \\
۲ & پیوند داده & پیوند \\
۱ & فیزیکی & پیوند \\
\خط‌پر
\پایان{جدول}

\برچسب{جدول:نگاشت‌لایه‌ها}
\پایان{لوح}

این نگاشت یک برابری کامل نیست، زیرا مرز مسئولیت‌ها در پیاده‌سازی‌های واقعی می‌تواند جابه‌جا شود.
\لر{OSI} هفت لایه دارد و طبق \لر{ISO/IEC 7498-1} یک مدل مرجع عمومی برای هماهنگ‌کردن استانداردهای ارتباط سیستم‌های باز است؛
خود استاندارد تصریح می‌کند که مدل، مشخصات یک پیاده‌سازی نیست [۷].
در مقابل، \لر{TCP/IP} هم معماری و هم مجموعه‌ای از پروتکل‌های عملی و مستقر در اینترنت است.
\لر{OSI} وظایف نمایش و نشست و همچنین پیوند داده و فیزیکی را جدا می‌کند، اما \لر{TCP/IP} آن‌ها را در لایه‌های کاربرد و پیوند ادغام می‌کند.
هر دو از لایه‌بندی، رابط سرویس و کپسوله‌سازی برای کاهش پیچیدگی و استقلال نسبی اجزا استفاده می‌کنند؛
بااین‌حال \لر{TCP/IP} از نظر مرزبندی عمل‌گراتر و کم‌لایه‌تر است.

\subsection{علت کارکرد کابل \لر{Straight} در شبکه‌های امروزی}

در اترنت ۱۰ و ۱۰۰ مگابیتی قدیمی، درگاه \لر{MDI} رایانه روی یک زوج ارسال و روی زوج دیگر دریافت می‌کرد،
درحالی‌که درگاه \لر{MDI-X} سوئیچ جای زوج‌های ارسال و دریافت را معکوس می‌کرد.
بنابراین اتصال دو تجهیز ناهم‌نوع، مانند رایانه و سوئیچ، با کابل \لر{Straight} و اتصال دو تجهیز هم‌نوع، مانند دو رایانه، با کابل \لر{Crossover} انجام می‌شد.

بیشتر \لر{PHY}های امروزی قابلیت \لر{Auto-MDI/MDIX} دارند.
این مدار هنگام برقراری پیوند تشخیص می‌دهد کابل مستقیم است یا متقاطع و در صورت لزوم نقش کانال‌های ارسال و دریافت را داخل تراشه جابه‌جا می‌کند.
برای نمونه، مستند رسمی \لر{PHY} مدل \لر{DP83826} شرکت \لر{Texas Instruments} بیان می‌کند که تراشه نوع کابل \لر{Straight} یا \لر{Crossover} را تشخیص می‌دهد و کانال‌های \لر{MDI/MDIX} را خودکار بازتخصیص می‌دهد؛
این سازوکار در بند \لر{40.8.2} استاندارد \لر{IEEE 802.3} توصیف شده است [۸].
در نتیجه، اگر دست‌کم یکی از دو سر پیوند \لر{Auto-MDI/MDIX} داشته باشد، هر دو نوع کابل معمولاً پیوند صحیحی ایجاد می‌کنند.

در \لر{1000BASE-T} نیز هر چهار زوج به‌طور هم‌زمان و دوطرفه استفاده می‌شوند و هر زوج با روش‌های جداسازی ارسال و دریافت، یک کانال تمام‌دوطرفه است [۳].
\لر{PHY}های گیگابیتی جدید علاوه بر مذاکره‌ی سرعت و حالت \لر{Duplex}، نگاشت زوج‌ها را نیز مدیریت می‌کنند؛
ازاین‌رو کابل \لر{Straight} به انتخاب متداول و عمومی تبدیل شده و دیگر لازم نیست کاربر برای بیشتر تجهیزات جدید کابل \لر{Crossover} جداگانه تهیه کند.

این قابلیت خرابی کابل را اصلاح نمی‌کند:
ترتیب دو سر کابل \لر{Straight} باید هر دو \لر{T568A} یا هر دو \لر{T568B} باشد، زوج‌ها نباید شکافته شوند و اتصال هر هشت رشته باید سالم باشد.
همچنین برخی تجهیزات قدیمی یا \لر{PHY}هایی که \لر{Auto-MDI/MDIX} در آن‌ها غیرفعال شده است هنوز به کابل \لر{Crossover} نیاز دارند.
پس علت اصلی کارکرد کابل‌های \لر{Straight}، تشخیص و جابه‌جایی خودکار \لر{MDI/MDIX} در لایه‌ی فیزیکی است، نه حذف مفهوم زوج‌های ارسال و دریافت.

\section{جمع‌بندی}

زوج‌به‌هم‌تابیده به‌دلیل هزینه‌ی کم و نصب ساده، انتخاب معمول برای دسترسی محلی تا ۱۰۰ متر است؛
فیبر برای برد، پهنای باند یا مصونیت بیشتر در برابر نویز برتری دارد؛
و کابل هم‌محور بیشتر در زیرساخت‌های موجود و کاربردهای تخصصی توجیه می‌شود.
مدل \لر{TCP/IP} وظایف ارتباطی را در چهار لایه‌ی عملی سازمان می‌دهد و نسبت به مدل مرجع هفت‌لایه‌ای \لر{OSI} ادغام بیشتری دارد.
همچنین فراگیرشدن \لر{Auto-MDI/MDIX} سبب شده است کابل مستقیم در بیشتر اتصال‌های اترنت جدید بدون توجه به نوع دو تجهیز کار کند.

\section{منابع}

\شروع{شمارش}

\فقره
\href{https://www.ieee802.org/misc-docs/GlobeCom2009/IEEE_802d3_Law.pdf}{\لر{IEEE 802.3, ``Ethernet: A Living Standard''}}، جدول تحول نرخ، فاصله و رسانه‌های اترنت.

\فقره
\href{https://www.belden.com/products/cable/coax-triax-cable/50-ohm-coax-cable/mrg5800}{\لر{Belden, MRG5800 RG-58 Coaxial Cable Specifications}}.

\فقره
\href{https://www.ieee802.org/3/minutes/july98/E2_0798.pdf}{\لر{IEEE 802.3ab, ``1000BASE-T Basics''}}، استفاده از چهار زوج، نرخ ۱ گیگابیت‌برثانیه و برد ۱۰۰ متر.

\فقره
\href{https://www.cisco.com/c/en/us/td/docs/sanity/suman/CPwE-Physical-Infrastructure_Feb2023.html}{\لر{Cisco, ``Physical Infrastructure Network Design for CPwE''}}، مقایسه‌ی برد، نرخ، نویز و هزینه‌ی مس و فیبر.

\فقره
\href{https://www.itu.int/rec/T-REC-G.652-202408-I/en}{\لر{ITU-T Recommendation G.652 (08/2024)}}، مشخصات فیبر و کابل نوری تک‌حالته.

\فقره
\href{https://www.rfc-editor.org/rfc/rfc1122.html}{\لر{IETF RFC 1122, ``Requirements for Internet Hosts --- Communication Layers''}}.

\فقره
\href{https://www.iso.org/standard/20269.html}{\لر{ISO/IEC 7498-1:1994, ``Open Systems Interconnection --- Basic Reference Model''}}.

\فقره
\href{https://www.ti.com/lit/ds/symlink/dp83826ai.pdf}{\لر{Texas Instruments, DP83826 Datasheet, ``Auto-MDIX Resolution''}}.

\پایان{شمارش}
